\section{From two adjacent intervals to gap counts}\label{sec:counting}

For an integer $X\ge2$ and a fixed $\lambda>0$, write
\[
 L=\log X,\qquad h=\lfloor\lambda L\rfloor,\qquad
 J=\{1,\ldots,h\},\qquad I=\{h+1,\ldots,2h\}.
\]
We use the uniform average
$\mathbb E_X A=X^{-1}\sum_{X<m\le2X}A(m)$.
Set $\vartheta(n)=(\log n)\mathbf1_{\{n\ {\rm prime}\}}$ and
\[
 V_B(m)=\frac1L\sum_{b\in B}\vartheta(m+b)\qquad(B=I,J).
\]
The condition $V_B>0$ says exactly that the corresponding interval
contains a prime. Moreover, $V_B\ge1$ whenever $V_B>0$, since all
its prime arguments exceed $X$.

\begin{proposition}[Weights for adjacent intervals]\label{prop:weights}
For each fixed $\lambda>0$ there is a finite constant $K_\lambda>0$
with the following property. For every $\varepsilon>0$, there are
weights $W_X(m)\ge0$, defined for integers $X<m\le2X$, such that
\begin{align}
 \mathbb E_X W_X&\longrightarrow1,
 &\mathbb E_X(W_XV_J)&\longrightarrow\lambda,\label{eq:weight-first}\\
 \limsup_{X\to\infty}\mathbb E_X(W_XV_I)&\le\varepsilon,
 &\limsup_{X\to\infty}\mathbb E_X(W_XV_J^2)&\le K_\lambda,\label{eq:weight-primes}\\
 \mathbb E_X W_X^2&=O_{\lambda,\varepsilon}(1).\label{eq:weight-second}
\end{align}
The constant $K_\lambda$ is independent of $\varepsilon$.
\end{proposition}

The construction and proof occupy Sections~\ref{sec:weights}
through~\ref{sec:moments}. We now show why precisely these estimates
suffice. In particular, the bound in \eqref{eq:weight-second} need
not be uniform as $\varepsilon$ decreases.

The estimate for $\mathbb E_X(W_XV_J^2)$ prevents the detected prime
mass from being concentrated on too little weighted mass. After removing
the starts that contain a prime in $I$, the separate estimate for
$\mathbb E_XW_X^2$ prevents the remaining weighted mass from being
concentrated on too few ordinary starts. This is why only the first
constant must be chosen before the suppression tolerance.

\begin{lemma}\label{lem:adjacent-gaps}
Suppose $h\ge1$ and that for at least $\delta X$ integers $m\in(X,2X]$, the interval
$(m,m+h]$ contains a prime and $(m+h,m+2h]$ contains none.
Then at least $\delta X/h$ distinct primes $p\in(X,2X+h]$ have a
consecutive successor $p^+$ satisfying $p^+-p>h$.
\end{lemma}
\begin{proof}
For each such $m$, let $p(m)$ be the last prime in $(m,m+h]$.
There is no prime in $(p(m),m+2h]$, so
$p(m)^+>m+2h\ge p(m)+h$.
For a fixed selected prime $p$, the condition $m<p\le m+h$ implies
$p-h\le m\le p-1$. This interval contains exactly $h$ integers.
Thus each prime is selected at most $h$ times, proving the claim.
\end{proof}

\begin{figure}[ht]
\centering
\setlength{\unitlength}{1pt}
\begin{picture}(420,91)
 \put(10,36){\line(1,0){390}}
 \put(10,32){\line(0,1){8}}
 \put(205,32){\line(0,1){8}}
 \put(400,32){\line(0,1){8}}
 \put(10,16){\makebox(0,0){$m$}}
 \put(205,16){\makebox(0,0){$m+h$}}
 \put(400,16){\makebox(0,0){$m+2h$}}
 \put(95,49){\makebox(0,0){contains a prime}}
 \put(305,49){\makebox(0,0){no primes}}
 \put(148,36){\circle*{5}}
 \put(148,23){\makebox(0,0){$p(m)$}}
 \put(148,67){\vector(1,0){267}}
 \put(283,81){\makebox(0,0){no next prime before $m+2h$}}
\end{picture}
\caption{The last prime in the first interval begins a gap longer than
$h$. Its distance from $m$ is at most $h$, which bounds the counting
multiplicity even when the next prime is much farther away.}
\label{fig:blocks}
\end{figure}

\begin{proof}[Proof of Theorem~\ref{thm:main}]
Fix $C>0$ and choose $\lambda>\max\{C,1\}$.
Take $\varepsilon=\lambda^2/(2K_\lambda)$ in
Proposition~\ref{prop:weights}, and fix the resulting weight family.
There is a finite constant $M>0$ such that
$\mathbb E_XW_X^2\le M$ for every sufficiently large $X$.
Let $A_X=\{V_J>0\}$ and $B_X=\{V_I>0\}$.
Weighted Cauchy--Schwarz gives
\[
 \mathbb E_X(W_X\mathbf1_{A_X})
 \ge \frac{\bigl(\mathbb E_X(W_XV_J)\bigr)^2}
             {\mathbb E_X(W_XV_J^2)}
\]
for all sufficiently large $X$; its denominator is positive because
the numerator tends to $\lambda^2>0$.
By \eqref{eq:weight-first}--\eqref{eq:weight-primes},
\[
 \liminf_{X\to\infty}\mathbb E_X(W_X\mathbf1_{A_X})
 \ge\frac{\lambda^2}{K_\lambda}.
\]
Since $\mathbf1_{B_X}\le V_I$ and $W_X\ge0$, it follows that
\[
 \liminf_{X\to\infty}\mathbb E_X
      (W_X\mathbf1_{A_X\setminus B_X})
 \ge \frac{\lambda^2}{K_\lambda}-\varepsilon
 =:\Delta>0.
\]
Another use of Cauchy--Schwarz yields, for every sufficiently large $X$,
\[
 \frac{|A_X\setminus B_X|}{X}
 \ge\frac{\bigl(\mathbb E_X
             (W_X\mathbf1_{A_X\setminus B_X})\bigr)^2}
           {\mathbb E_XW_X^2}
 \ge\frac{\Delta^2}{4M}=:\delta>0.
\]
This is an ordinary proportion of integer starting points.
Lemma~\ref{lem:adjacent-gaps} therefore gives at least
\[
 \frac{\delta X}{h}\ge\frac{\delta X}{\lambda\log X}
\]
distinct prime starts $p\in(X,2X+h]$ with consecutive gap exceeding $h$.
Uniformly on this range,
$\log p=\log X+O(1)$ and $h=\lambda\log X+O(1)$.
The strict inequality $\lambda>C$ consequently gives
$h>C\log p$ for every sufficiently large $X$.

For an arbitrary sufficiently large integer $N$, choose
$X=\lfloor p_N/3\rfloor$. Then $2X+h<p_N$, so all these prime starts
have indices at most $N$. The prime number theorem gives
\[
 \frac{X}{\log X}\sim\frac{p_N}{3\log p_N}\sim\frac N3.
\]
In particular $X/\log X\ge N/4$ for every sufficiently large $N$.
The required count is at least $\delta N/(4\lambda)$.
Thus $c(C)=\delta/(4\lambda)>0$ is admissible.
All choices preceding $X$ depend only on the fixed threshold.
\end{proof}

The rest of the paper proves Proposition~\ref{prop:weights}.
The task is to make the weighted prime mass in $I$ arbitrarily small
without losing either the mass of the square or the uniform detector
bound in the adjacent block $J$.
